\documentclass[10pt,a4paper]{amsart}
\usepackage[margin=19mm]{geometry}
\usepackage{amsmath,amssymb,mathtools}
\usepackage[scr=rsfs,cal=euler]{mathalfa}
\usepackage{xfrac}
\usepackage[unicode,hidelinks,bookmarks=false]{hyperref}
\newcommand{\ie}{i.e.\ }
\newcommand{\cf}{cf.\ }
\newcommand{\resp}{respectively\ }
\newcommand{\Subsection}{\subsection}
\providecommand{\nobreakdash}{\nobreak\hbox{-}}
\setlength{\emergencystretch}{2em}
\multlinegap0pt
\usepackage[shortlabels]{enumitem}
\usepackage{enumerate}
\usepackage{natbib}
\usepackage{amssymb}
\usepackage{mathtools}
\usepackage{bbm}
\newtheorem{lemma}{Lemma}
\usepackage{cleveref}
\crefname{lemma}{Lemma}{Lemmas}
\newcommand{\R}{\mathbb{R}}
\newcommand{\abs}[1]{\left|#1\right|} % absolute value
\newcommand{\cH}{\mathcal{H}}
\newcommand{\norm}[1]{\left\|#1 \right\|} % Norm
\def\red{\color{red}}
\def\ud{\mathrm d}
\title{Convergence of the generalization error for Deep Gradient Flow Methods for PDE\MakeLowercase{s}}
\author{Chenguang Liu, Antonis Papapantoleon, Jasper Rou}
\date{Source: arXiv:2512.25017v2 (CC BY 4.0)}
\begin{document}
\maketitle

\noindent\emph{Source and licence.} Convergence of the generalization error for Deep Gradient Flow Methods for PDE\MakeLowercase{s}. Version arXiv:2512.25017v2 (CC BY 4.0), distributed by arXiv under the Creative Commons Attribution 4.0 International licence, http://creativecommons.org/licenses/by/4.0/, as recorded in the arXiv licence metadata for this exact version and shown on the abstract page https://arxiv.org/abs/2512.25017v2. Full licence notice: \texttt{licenses/arxiv-2512.25017-CC-BY-4.0.txt}.\par\smallskip

\noindent\textit{Editorial scope.} Counted unit: the article's lemma lem:lip\_Xtheta (source0.tex lines 2026--2032). The article's complete original proof of it is delivered with it (source0.tex lines 2034--2106, one \texttt{proof} environment of 73 lines). No mathematics is written, repaired or re-derived: the statement and the proof are the article's own lines, in the article's own notation, and no retained line is substituted or re-flowed.  Retained uncounted as the source-local context the proof needs: lines 1858--1867.  Nothing was omitted from any retained span, so the package carries no omission record.  No retained line contains a citation, so the wrapper carries no bibliography and no bibliography entry.  No retained line contains a figure environment, an image-inclusion command or drawing code, so no image is delivered and the package declares no asset dependency.  Editorial formatting only, in addition: an AMS \texttt{amsart} preamble that restates the article's own declarations and macros used by the retained lines, namely the environments \texttt{lemma} on separate counters, together with the standard packages those lines need.  The retained environments print in this document as Lemma 1 in order of appearance, whereas the article prints the same items with the numbering of its own full document.  The complete native source of the article as distributed in the arXiv e-print for this exact version is bundled unchanged as \texttt{sources/191958/source0.tex}.
\begin{lemma}[$\cH_0^1$-boundedness of $X^n$]
\label{lem:boundXN}
Let $\theta \in \R \times \R \times \R^d$, then $X^n$ is bounded in $\cH_0^1$, \emph{i.e.} there exists a constant $C_\psi>0$ such that
\begin{align*}
% \red \abs{X^n(\theta)} \le 
% &\red \abs{\psi \left( \hat\alpha x +  \hat c \right)} + r_n \abs{\left( x\cdot\nabla \psi \right) \left( \hat\alpha x + \hat c \right)} + r_n \abs{\left( \nabla \psi \right) \left( \hat\alpha x + \hat c \right)}\le C_\psi (r_n)^3, \\
% \intertext{and }
\norm{X^n(\theta)}^2_{\cH_0^1} &\le C_\psi \left(r_n\right)^{d+8}.
\end{align*}
\end{lemma}

\begin{lemma}[$\theta$-Lipschitz continuity]
\label{lem:lip_Xtheta}
Let $\theta,\theta'\in \R\times\R\times\R^d$, then $X^n$ is Lipschitz continuous in $\cH_0^1$, \emph{i.e.} there exists a constant $C_\psi>0$ such that
\begin{align*}
    \norm{X^n(\theta) - X^n(\theta')}_{\cH_0^1} \le C_\psi \left(r_n\right)^{4+\frac{d}{2}}\abs{\theta-\theta'}^{\frac{1}{2}}.
\end{align*}
\end{lemma}

\begin{proof}
% The derivatives with respect to the neural network parameters are
% \[
% \begin{aligned}
% \frac{\partial}{\partial \beta} \beta \psi \left( \alpha . + c \right) & = \psi \left( \alpha x + c \right), \\ 
% \frac{\partial}{\partial \alpha} \beta \psi \left( \alpha . + c \right) & = \beta x^\mathsf{T} \left( \nabla \psi \right) \left( \alpha x + c \right), \\ 
% \frac{\partial}{\partial c} \beta \psi \left( \alpha . + c \right) & =  \beta \left( \nabla \psi \right) \left( \alpha x + c \right).
% \end{aligned}
% \]
% The parameter derivatives of $\nabla_x \psi \left( \alpha x + c \right)$ are
% \[
% \begin{aligned}
% \frac{\partial}{\partial \beta} \nabla_x \left(\beta \psi \left( \alpha . + c \right)\right) & = \alpha \left( \nabla \psi \right) \left( \alpha x + c \right), \\ 
% \frac{\partial}{\partial \alpha}  \nabla_x \left(\beta \psi \left( \alpha . + c \right)\right) & = \beta \alpha x^\mathsf{T} D^2 \psi \left( \alpha x + c \right) {\color{orange} + \beta \left( \nabla \psi \right) \left( \alpha x + c \right)}, \\ 
% \frac{\partial}{\partial c} \nabla_x \left(\beta \psi \left( \alpha . + c \right)\right) & =\beta\alpha D^2 \psi \left( \alpha x + c \right). \\ 
% \end{aligned}
% \]
As in the proof of \cref{lem:boundXN}, we have
\begin{align}
\label{eq:appendix-bound}
\norm{X^n(\theta) - X^n\left( \theta'\right)}_{\cH_0^1} 
    &\le \norm{\psi \left( \hat\alpha \cdot + \hat c \right)-\psi \left( \hat\alpha' \cdot + \hat c'\right)}_{\cH_0^1} \nonumber \\ \nonumber
    &\qquad  +\norm{\hat\beta  \left( \cdot\nabla \psi \right) \left( \hat\alpha \cdot + \hat c \right)-\hat\beta' \cdot \left( \nabla \psi \right) \left( \hat\alpha' \cdot + \hat c' \right)}_{\cH_0^1} \\
    &\qquad + \norm{\hat\beta \left( \nabla \psi \right) \left( \hat\alpha \cdot + \hat c \right)-\hat\beta' \left( \nabla \psi \right) \left( \hat\alpha' \cdot + \hat c' \right)}_{\cH_0^1}.
\end{align}
Let us recall that $\abs{\hat \alpha}, |\hat\beta|,\abs{\hat c}\le r_n$ and $\abs{\hat \alpha}\ge r_n^{-1}$. 
Using that $\psi \in C_c^{\infty} \left( \mathbb{R}^d \right)$, and therefore Lipschitz, we get
\begin{align*}
\int_{\R^d} \abs{\psi \left( \hat\alpha x + \hat c \right)-\psi \left( \hat\alpha' x + \hat c'\right)}^2\ud x 
    & \le \int_{\R^d} C_\psi \abs{\left(\hat\alpha - \hat \alpha' \right) x + \hat c - \hat c'} \abs{\psi \left( \hat\alpha x + \hat c \right) - \psi \left( \hat\alpha' x + \hat c'\right)} \ud x \\
    & \le C_\psi \abs{\theta-\theta'} \int_{\R^d} \left( 1 + \abs{x} \right) \left( \abs{\psi \left( \hat\alpha x + \hat c \right)}+ \abs{\psi \left( \hat\alpha' x + \hat c'\right)} \right) \ud x.
\end{align*}
Using the change of variables $y = \hat \alpha x + \hat c$, we arrive at
\begin{align*}
    \int_{\R^d} \left( 1 + \abs{x} \right) \abs{\psi \left( \hat\alpha x + \hat c \right)} \ud x & = \abs{\hat \alpha}^{-d} \int_{\R^d} \abs{\psi (y)} \ud y + \abs{\hat \alpha}^{-d-1} \int_{\R^d} \abs{y - \hat c} \abs{\psi \left( y \right)} dy \\
    & \le C_\psi \Big( r_n^{d} + r_n^{d+2} \Big),
\end{align*}
and we obtain the bound
\[
    \int_{\R^d} \abs{\psi \left( \hat\alpha x + \hat c \right)-\psi \left( \hat\alpha' x + \hat c'\right)}^2\ud x \leq C_\psi \abs{\theta-\theta'} r_n^{d+2}.
\]
Analogously, using that $\nabla \psi$ is Lipschitz as well, we have that
\begin{multline*}
\int_{\R^d} \abs{\hat \alpha\nabla\psi\left( \hat\alpha x + \hat c \right)-\hat \alpha'\nabla\psi \left( \hat\alpha' x + \hat c'\right)}^2\ud x \\
     \le 2\int_{\R^d} \abs{\hat\alpha-\hat \alpha'}^2\abs{\nabla\psi \left( \hat\alpha x + \hat c \right)}^2\ud x + 2\int_{\R^d}\abs{\hat \alpha'}^2\abs{\nabla\psi \left( \hat\alpha x + \hat c \right)-\nabla\psi \left( \hat\alpha' x + \hat c'\right)}^2\ud x \\
      \le 4 r_n \abs{\theta - \theta'} \int_{\R^d} \abs{\nabla\psi \left( \hat\alpha x + \hat c \right)}^2\ud x + 2 r_n^2 \int_{\R^d} \abs{\nabla\psi \left( \hat\alpha x + \hat c \right)-\nabla\psi \left( \hat\alpha' x + \hat c'\right)}^2\ud x \\
     \le  C_\psi \abs{\theta-\theta'}(r_n)^{d+4}.
\end{multline*}
Therefore, we arrive at the following bound for the $\cH_0^1$-norm of this term
\[
    \norm{\psi \left( \hat\alpha \cdot + \hat c \right)-\psi \left( \hat\alpha' \cdot + \hat c'\right)}_{\cH_0^1}\le C_\psi \abs{\theta-\theta'}^{1/2}(r_n)^{d/2+2}.
\]
We can similarly estimate the other two terms in \eqref{eq:appendix-bound}, and we get that
\begin{align*}
&\norm{\hat\beta  \left( \cdot\nabla \psi \right) \left( \hat\alpha \cdot + \hat c \right)-\hat\beta' \cdot \left( \nabla \psi \right) \left( \hat\alpha' \cdot + \hat c' \right)}_{\cH_0^1} \\
    &\qquad + \norm{\hat\beta \left( \nabla \psi \right) \left( \hat\alpha \cdot + \hat c \right)-\hat\beta' \left( \nabla \psi \right) \left( \hat\alpha' \cdot + \hat c' \right)}_{\cH_0^1} 
    \le  C_\psi \abs{\theta-\theta'}^{1/2}(r_n)^{d/2+4}
\end{align*}
%Therefore we can bound the difference
 %   \begin{align*}
  %      \abs{X(\theta) - X \left( \widetilde \Theta \right)} \le & C_\psi \left( \abs{\alpha} + \abs{\tilde \alpha} + \abs{\beta} + \abs{\tilde \beta} \right) \left( \abs{x}^2 + 1 \right) \abs{\theta - \widetilde\Theta}, \\
    %    \abs{\nabla_x X(\theta) - \nabla_x X \left( \widetilde \Theta \right)} \le & C_\psi \left( \abs{\alpha}^2 + \abs{\tilde \alpha}^2 + \abs{\beta} + \abs{\tilde \beta} + 1 \right) \left( \abs{x}^2 + 1 \right) \abs{\theta - \widetilde\Theta}.
    %\end{align*}
    %As $\abs{\alpha},\abs{\tilde \alpha},\abs{\beta},\abs{\tilde \beta}$ are bounded by $r_n$ and $x$ is supported on $D^n$, 
     %\begin{align*}
        %\abs{X^n(\theta)-X^n\left(\widetilde \Theta\right)}\le& C_\psi\left(r_n\right)^{1+\frac{2}{d}}\abs{\theta-\widetilde\Theta}, \\
         % \abs{\nabla_x X(\theta)-\nabla_x X\left(\widetilde \Theta\right)}\le&C_\psi\left(r_n\right)^{2+\frac{2}{d}}\abs{\theta-\widetilde\Theta}.
    %\end{align*}
Concluding, we have
\[
    \norm{X^n(\theta) - X^n\left( \theta'\right)}_{\cH_0^1} \le C_\psi \left(r_n\right)^{4+\frac{d}{2}}\abs{\theta-\theta'}^{\frac{1}{2}}. \qedhere
\]
\end{proof}

\end{document}
